\documentclass[aspectratio=169,12pt]{beamer}
%	\usetheme{Antibes}
	\usetheme{Frankfurt}

%\usepackage{amsmath}

\usepackage{graphics}
\graphicspath{{../img}}

%\definecolor{VerdeApostila}{rgb}{0.058823529411764705, 0.6862745098039216, 0.30980392156862746}
%\usecolortheme[named=VerdeApostila]{structure}

\newcommand{\modo}[2]{
%	#2
} % claro ou escuro
\newcommand{\quebra}{
	\break
}

\newcommand{\iniciar}[3]{

	\begin{frame}[plain]
		\maketitle
		\centering
		
		\textbf{Aula #1}: 
		#2 
		 
		#3
	  \end{frame}

%	\section*{Nesta aula}
	\begin{frame}{Nesta aula}
		\begin{multicols}{3}
			\tableofcontents
		\end{multicols}	
	  \end{frame}  

}

\newcommand{\terminar}[1]{

\section{Exercícios #1}
\begin{frame}[allowframebreaks]{Exercícios}	% containsverbatim
	\input{aula#1_exercícios}	
  \end{frame}	
  
\section*{}  

\begin{frame}[plain]
	\maketitle
	
	\centering	
	\Large
	\textbf{Dúvidas?}
		
	\url{akira.demenech@uel.br}
  \end{frame}

\end{document}

}

%\usepackage[brazil]{babel}	
%\usepackage[T1]{fontenc}

\usepackage{xcolor}

\usepackage{amsmath}
\usepackage{mathdots}

\usepackage{multicol}

\usepackage{hyperref}
\newcommand{\link}[2]{\href{#1}{\texttt{#2}}}

\title{Minicurso de Python}
\author{Guilherme Akira Demenech Mori}
\date{}

\institute{
%	\link{https://www.facebook.com/codificajr/}{
%		\includegraphics[width=0.3\textwidth]{\modo{Logo pronta.pdf}{logo clara.pdf}}
%	}
}

\begin{document}
	\modo{}{\color{white}}

\iniciar{2}{Condicionais e Laços}{Fluxo}

\section{Comparações}
\begin{frame}{Mais operações}
	
	Para valores de todos os tipos:
	
%	\begin{multicols}{2}
		\begin{itemize}
			\item \textbf{Igualdade}: 
				\texttt{x == y} 
			
			\item \textbf{Desigualdade}: 
				\texttt{x != y}
			
		\end{itemize}	
%	\end{multicols}	
	
	\vfill
	
	Para valores \textbf{int}, \textbf{float}, \textbf{str}, \textbf{tuple}, \textbf{list}:
	
	\begin{multicols}{2}
		\begin{itemize}
			\item \textbf{Menor que}: 
				\texttt{x < y} 
			
			\item \textbf{Menor ou igual que}: 
				\texttt{x <= y}
			
			\item \textbf{Maior que}: 
				\texttt{x > y} 
			
			\item \textbf{Maior ou igual que}: 
				\texttt{x >= y}									
		\end{itemize}	
	\end{multicols}			
  \end{frame}
  
\begin{frame}{Contém} 
	Verificar se o composto (\textbf{tuple}, \textbf{list} ou \textbf{set}) contém o valor:
	
	\texttt{valor in nomevariavel}
	
\vfill	
	
	Verificar se o \textbf{dict} contém a chave:
	
	\texttt{chave in nomevariavel}
	
\vfill	
	
	Verificar se uma \textbf{str} contém uma outra \textbf{str}:
	
	\texttt{valors in nomevariavels}
  \end{frame}	

\section{Condicionais} 
\begin{frame}[containsverbatim]{Se .... senão ...}
%	\begin{multicols}{2}			
	
	Executa as linhas dentro do bloco \texttt{if} se a condição for satisfeita.	
	
	\begin{verbatim}
if condicao:
		acaosesim
	\end{verbatim}
	
	\vfill

	Executa as linhas dentro do bloco \texttt{else} se a condição não for satisfeita.	

\begin{verbatim}
if condicao:
		acaosesim
else:	
		acaosenao
\end{verbatim}	
%	\end{multicols}	
  \end{frame}
  
\begin{frame}[containsverbatim]{Condicionais aninhados}  			
	
%	\begin{multicols}{2}
	Executa as linhas dentro do bloco \texttt{elif} se as condições anteriores não forem satisfeitas, mas a sua for.	
	
	\texttt{else} é executada se nem \texttt{if}, nem nenhum \texttt{elif} forem.
	
	\begin{verbatim}
if condicao1:
		acaosesim1
elif condicao2:	
		acaosenao1sim2
elif condicao3:
		acaosenao12sim3	
else:	
		acaosenao123
	\end{verbatim}
\end{frame}	
\begin{frame}[containsverbatim]%{Condicionais aninhados}	

	Você pode colocar condicionais dentro de condicionais também.
	
	\begin{verbatim}
if condicao1:
		if condicao2:
				acaosesim12
		else:	
				acaosesim1nao2
else:	
		if condicao3: 
				acaosenao1sim3
		else:	
				acaosenao13
	\end{verbatim}
%	\end{multicols}	
  \end{frame}

\section{Laços de repetição}

\begin{frame}[containsverbatim]{Enquanto}				
		
		Executa as linhas dentro do bloco \texttt{while} enquanto a condição for satisfeita.	
		
		\begin{verbatim}
while condicao:
		acaoenquantosim
		\end{verbatim}
		
		Lembre-se de programar seu código de maneira que laços \texttt{while} não fiquem presos sem alteração da condição.		
				 				 		
  \end{frame}

\begin{frame}[containsverbatim]{Para ... em ....}				
	
	Executa as linhas dentro do bloco \texttt{for} para cada valor de um composto, atribuindo a uma variável.	
	
	\begin{verbatim}
for nomevariavel in composto:
		acaoparacadavalor
	\end{verbatim}				
	
%	\vfill 
	
	
	
  \end{frame}
  
\frame{
	A função \texttt{range(n)} cria um composto com os inteiros de \texttt{0} até \texttt{n-1}, nessa ordem, e pode ser utilizado no \texttt{for}.
	\hfill 
	\# 
	igual a \texttt{range(0,n)} e \texttt{range(0,n,1)}
	
	
	\texttt{range(m,n)} com os inteiros de \texttt{m} até \texttt{n-1}.
	\hfill 
	\# 
	igual a \texttt{range(m,n,1)}
	
	\texttt{range(m,n,s)} com os inteiros entre \texttt{m} e \texttt{n-1} começando em \texttt{m} e somando \texttt{s}.
}  
  
\begin{frame}[containsverbatim]{Quebras}				
\begin{multicols}{2}	
	
	\texttt{break} sai da repetição e vai para a linha depois do fim do laço	
	
\begin{verbatim}
while condicao:
		acaoenquantosim
		
		if condicaoparada:
				break 
				
		maisacaoenquantosim
		
acoesposteriores				
\end{verbatim}	
	
	\texttt{continue} pula o resto do bloco e vai para a próxima repetição, do começo do laço 				

\begin{verbatim}
for nomevariavel in composto:
		acaoparacadavalor
	
		if condicaodesalto:
				continue 
		
		maisacaoparacadavalor						
\end{verbatim}

\end{multicols}	
\end{frame}  

% Tabs ou espaços?

\section{Indentação}
\begin{frame}{Tabs ou espaços?}
	
	\Large
	Para avançar um nível (um condicional ou laço, por exemplo), dê tab ou 4 espaços.
	
	\LARGE
	Utilize consistentemente, ou só tabs, ou só 4 espaços.
	
%	\vfill 
	
	
  \end{frame}
  
\frame{\Huge Aninhe estruturas condicionais e laços à vontade.}  


\terminar{2}